\providecommand{\titlepagecoverimg}{img/auxl/cover.png}

\providecommand{\prefacenotesIO}{}
\providecommand{\prefacenotesinteractive}{}
\providecommand{\prefacenotesadditional}{}

% add this as \prefacenotesIO for finals
% We recommend learning and using these functions during the Practice Session.

% possible item to add, \prefacenotesadditional
% \item All problems are solvable in C++, but the same is not guaranteed for Python due to its slowness.



%%%%%%%%%%%%%%%%%%%%%%%%%%%%%%%%%%%%%%%%%%%%%%%%%%%%%%%%%
% title page
%%%%%%%%%%%%%%%%%%%%%%%%%%%%%%%%%%%%%%%%%%%%%%%%%%%%%%%%%

\begin{titlepage}
\vspace*{4.65in}
\setBG{\includegraphics[width=\pagewidth]{\titlepagecoverimg}}
\end{titlepage}




%%%%%%%%%%%%%%%%%%%%%%%%%%%%%%%%%%%%%%%%%%%%%%%%%%%%%%%%%
% long notes
%%%%%%%%%%%%%%%%%%%%%%%%%%%%%%%%%%%%%%%%%%%%%%%%%%%%%%%%%

% \newcommand\preem{rage}
\newcommand\preem{smiling-face-with-sunglasses}

\cleardoublepage

\markboth{Important!}{}
% \vspace*{0.05in}
% {\Large \color{red}{\textbf{Important!} Read the following:}}
% {\Huge {\color{red}{\textbf{Important!} Read the following: {\fontsize{80}{96}\texttt{\emoji{\preem}}} }} }
% {\Huge {\color{red}{\textbf{{\texttt{\emoji{rage}} Important!}} Read the following: }} }

% \newlength{\ragewidth}
% \newlength{\rageheight}
% \settowidth{\ragewidth}{\scalebox{3.3}{\texttt{\emoji{\preem}}}}
% \settoheight{\rageheight}{\texttt{\emoji{\preem}}}
% {\Huge{\color{red}{\textbf{%
%     % \makebox[\ragewidth][r]{\scalebox{2}{\texttt{\emoji{\preem}}}}
%     \parbox[b][\rageheight]{\ragewidth}{
%         \makebox[\ragewidth][r]{
%         \lapbox[\width]{17pt}{\raisebox{2pt}{\scalebox{5.2}{\texttt{\emoji{\preem}}}}}}}
%     Important!} Read the following:}}}

\textbf{Access the contest website} at the following url:
\begin{center}
    \url{https://tinyurl.com/2025AlgolympicsFinalsPractice}
\end{center}

\begin{description}

\item[Practice Round] This is the time to test the contest environment.
\begin{itemize}
    \item Familiarize yourself with the text editor, and practice compiling/running locally.
    \item Practice submissions and using the clarification system on DOMJudge.
\end{itemize}

%\item[Test Cases.]
\item[Hidden Test Cases.]
Your solution will be checked by running it against one or more (usually several) hidden test cases. You will not have access to these cases, but a correct solution is expected to handle them correctly.

% \item[Output Format.]
\item[Strict Output Format.]
The output checker is {\color{red}{\textbf{strict}}}. Follow these guidelines strictly:

\begin{itemize}
    \item It is \textbf{space sensitive}. Do not output extra leading or trailing spaces. Do not output extra blank lines unless explicitly stated.
    \item It is \textbf{case sensitive}. So, for example, if the problem asks for the output in lowercase, follow it.
    \item Ensure that your output ends with a newline (by ending with \verb|cout << endl| perhaps, in C++; this is already the default behavior of \verb|print| in Python)
    \item Do not print any tabs. (No tabs will be required in the output.)
    \item Do not output anything else aside from what's asked for in the Output section. So, do not print things like ``\texttt{Please enter t}''. 
\end{itemize}

Not following the output format strictly and exactly will likely result in the verdict ``\textit{Output isn't correct}''.

% \item[Standard I/O.]
\item[Use Standard I/O.]
Do not read from, or write to, a file. You must read from the standard input and write to the standard output.

\item[Submit Code Only.]
Only include \textbf{one} file when submitting: the source code (.cpp, .py, etc.) and nothing else.

\item[No Java Package.]
For Java submissions, do not include a \java{package} line.

% \item[Only Builtin Libraries.]
% Do not use third-party libraries. Only libraries that come installed with the compilers are available on the evaluation machines.

% \item[Filename.]
\item[No Weird Filenames.]
Only use letters, digits and underscores in your filename. Do not use spaces or other special symbols.

% \item[Fast I/O.]
\item[Use Fast I/O.]
Many problems have large input file sizes, so use fast I/O. For example:
\begin{itemize}
\item In C/C++, use \cpp{scanf} and \cpp{printf}.
\item In Python, use \python{sys.stdin.readline()}
\end{itemize}

\prefacenotesIO{}

% \item[Interactive Problems.]
% \item[Flush On Interactive Problems.]
% On interactive problems, make sure to \textbf{flush} your output stream after printing.
% \begin{itemize}
%     % weird hack with mintinline since it's converting << to a special character. feel free to use better fix
%     \item In C++, use \cpp{fflush(stdout);} or \mintinline[escapeinside=||]{cpp}{cout <||< endl;}
%     \item In Python, use \python{sys.stdout.flush()} or \python{print(flush=True)}
%     \item For more details, including for other languages, ask a question/clarification through CMS.
% \end{itemize}

% \prefacenotesinteractive{}

% \prefacenotesadditional{}

\end{description}

Good luck and enjoy the contest!  \texttt{\emoji{smiley}}


% TODO
% "how compprog works"
%     link to absolute beginner tutorial (video?), 
%     link to absolute beginner tutorial in abakoda?
%     link to absolute beginner tutorial document of our own making
% link to more specific/obscure/technical stuff (like java package name)
% link to verdicts page ("Documentation" in CMS)







%%%%%%%%%%%%%%%%%%%%%%%%%%%%%%%%%%%%%%%%%%%%%%%%%%%%%%%%%
% table of contents, and short notes
%%%%%%%%%%%%%%%%%%%%%%%%%%%%%%%%%%%%%%%%%%%%%%%%%%%%%%%%%

\setBG{\contestbg}

% \textsf{\tableofcontents}
% \markboth{Table of Contents}{}

% \begin{comment}
% \section*{Notes}

% these notes are intentionally repeated from the "Very Important" section below.

% \begin{itemize}
% \item Many problems have large input file sizes, so use fast I/O. For example:
% \begin{itemize}
%     \item In C/C++, use \cpp{scanf} and \cpp{printf}.
%     \item In Python, use \python{sys.stdin.readline()}
% \end{itemize}

% \prefacenotesIO{}

% \item On interactive problems, make sure to \textbf{flush} your output stream after printing.
% \begin{itemize}
%     % weird hack with mintinline since it's converting << to a special character. feel free to use better fix
%     \item In C++, use \cpp{fflush(stdout);} or \mintinline[escapeinside=||]{cpp}{cout <||< endl;}
%     \item In Python, use \python{sys.stdout.flush()} or \python{print(flush=True)}
%     \item For more details, including for other languages, ask a question/clarification through CMS.
% \end{itemize}

% \prefacenotesinteractive{}

% \prefacenotesadditional{}
% \end{itemize}

% Good luck and enjoy the problems!
% \end{comment}
